\underline{Exercise 5.2}

% Question 1
\begin{enumerate}
\item 
\begin{enumerate}
\item[(c)] \[
\begin{aligned}
& 8\left|\begin{array}{ll}
0 & 1 \\
0 & 3
\end{array}\right|-1\left|\begin{array}{ll}
4 & 1 \\
6 & 3
\end{array}\right|+3\left|\begin{array}{ll}
4 & 0 \\
6 & 0
\end{array}\right| \\
=& 8(0-0)-1(12-6)+3(0-0) \\
=& 0-6+0 \\
= &-6 \\
\end{aligned}
\]

\item[(e)] \[
\begin{aligned}
& a\left|\begin{array}{ll}
c & a \\
a & b
\end{array}\right|-b\left|\begin{array}{cc}
b & a \\
c & b
\end{array}\right|+c\left|\begin{array}{cc}
b & c \\
c & a
\end{array}\right| \\
=& a\left(c b-a^{2}\right)-b\left(b^{2}-a c\right)+c\left(a b-c^{2}\right) \\
=& a b c-a^{3}-b^{3}+a b c+a b c-c^{3} \\
=&-a^{3}-b^{3}-c^{3}+3 a b c
\end{aligned}
\]

\item[(f)] \[
\begin{aligned}
& x\left|\begin{array}{cc}
y & 2 \\
-1 & 8
\end{array}\right|-5\left|\begin{array}{cc}
3 & 2 \\
9 & 8
\end{array}\right|+0\left|\begin{array}{rr}
3 & y \\
9 & -1
\end{array}\right| \\
=& x(8 y+2)-5(24-18)+0(-3-9 y) \\
=& 8 x y+2 x-30+0 \\
=& 8 x y+2 x-30
\end{aligned}
\]
\end{enumerate}

% Question 2
  \item \( \left|C_{13}\right|: 1+3=4 \) is even so \( + \)\\
\( \left|C_{23}\right|: 2+3=5 \) is odd so \(- \)\\
\( \left|C_{33}\right|: 3+3=6 \) is even so \( + \)\\
\( \left|C_{41}\right|: 4+1=5 \) is odd so \( - \)\\
\( \left|C_{34}\right|: 3+4=7 \) is odd so \( - \) \\


% Question 3
\item 
Minor of \(a\):
\[ \left|M_{11}\right|=\left|\begin{array}{cc}e & f \\ h & i\end{array}\right|=e i-f h \]
Cofactor of \(a\): \[
\left|C_{11}\right|=(-1)^{1+1}\left|M_{11}\right|=\left|M_{11}\right|
\]

Minor of \(b\):
\[ \left|M_{12}\right|=\left|\begin{array}{ll}d & f \\ g & i\end{array}\right|=d i-f g \]
Cofactor of \(b\):
\[\left|C_{12}\right|=(-1)^{3}\left|M_{12}\right|=-\left|M_{12}\right| \]
Minor of \(f\):
\[ \left|M_{23}\right|=\left|\begin{array}{ll}a & b \\ g & h\end{array}\right|=a h-b g \]
Cofactor of \(f\):
\[\left|C_{23}\right|=(-1)^{5}\left|M_{23}\right|=-\left|M_{23}\right| \] \\

% Question 6
\item[6.] Minors of third row:
\[
\begin{aligned}
\left|M_{31}\right| &=\left|\begin{array}{cc}
11 & 4 \\
2 & 7
\end{array}\right| = 69, \quad
\left|M_{32}\right|&=\left|\begin{array}{cc}
9 & 4 \\
3 & 7
\end{array}\right| = 51, \quad
 \left|M_{33}\right|&=\left|\begin{array}{ll}
9 & 11 \\
3 & 2
\end{array}\right| =-15
\end{aligned}
\]
Cofactors: $ |C_{31}| =  |M_{31}|, |C_{32}| =  -|M_{32}| , |C_{33}| =  |M_{33}| $.\end{enumerate}

%%%%%%%%%%%%%%%% Exercise 5.3

\underline{Exercise 5.3}

% Question 1
\begin{enumerate}
\item 
\[
\begin{aligned}
\left|\begin{array}{ccc}
4 & 0 & -1 \\
2 & 1 & -7 \\
3 & 3 & 9
\end{array}\right| &=4\left|\begin{array}{rr}
1 & -7 \\
3 & 9
\end{array}\right|-0\left|\begin{array}{cc}
2 & -7 \\
3 & 9
\end{array}\right|-1\left|\begin{array}{ll}
2 & 1 \\
3 & 3
\end{array}\right| \\
&=4(9+21)-0(18+21)-1(6-3) \\
&=120-0-3=117
\end{aligned}
\]

Interchanging rows and columns:\\

\( \begin{aligned}\left|\begin{array}{ccc}4 & 2 & 3 \\ 0 & 1 & 3 \\ -1 & -7 & 9\end{array}\right| &=4\left|\begin{array}{cc}1 & 3 \\ -7 & 9\end{array}\right|-2\left|\begin{array}{cc}0 & 3 \\ -1 & 9\end{array}\right|+3\left|\begin{array}{cc}0 & 1 \\ -1 & -7\end{array}\right| \\ &=4(9+21)-2(0+3)+3(0+1) \\ &=120-6+3=117 \end{aligned} \) \\

Interchange row 1 and 2:
\[
\begin{aligned}
\left|\begin{array}{rrr}
2 & 1 & -7 \\
4 & 0 & -1 \\
3 & 3 & 9
\end{array}\right| &=2\left|\begin{array}{rr}
0 & -1 \\
3 & 9
\end{array}\right|-1\left|\begin{array}{cc}
4 & -1 \\
3 & 9
\end{array}\right|-7\left|\begin{array}{ll}
4 & 0 \\
3 & 3
\end{array}\right| \\
&=2(0+3)-1(36+3)-7(12-0) \\
&=6-39-84=-117
\end{aligned}
\]
You can verify the other two properties in the same way. \vspace{1em}

% Question 4
\item[4.] Multiplying every element of an $n$th-order determinant by $k$ multiplies each of its $n$ rows by $k$. Multiplying a single row by $k$ multiplies the determinant by $k$, so doing it to all $n$ rows gives $k^{n}|A|$.

% Question 5
\item[5.] 
\begin{enumerate}
\item[(a)] \[
\begin{array}{l}
4\left|\begin{array}{cc}
1 & -3 \\
1 & 0
\end{array}\right|+1\left|\begin{array}{cc}
19 & 1 \\
7 & 1
\end{array}\right| \\
=4(0+3)+1(19-7) \\
=\quad 12+12=24
\end{array}
\]

Non-singular, rank = 3. \\

\item[(b)] \[
\begin{aligned}
& 5\left|\begin{array}{cc}
-2 & 1 \\
0 & 3
\end{array}\right|+6 \left| \begin{array}{ll}
4 & 1 \\
7 & 3
\end{array}\right| \\
=& 5(-6-0)+6(12-7) \\
=&-30+30=0
\end{aligned}
\]

Singular. Rank<3, will need to put in echelon form to find exact rank. \\

\item[(c)]
\[
\begin{aligned} & 7\left|\begin{array}{cc}1 & 4 \\ -3 & -4\end{array}\right|+1\left|\begin{array}{cc}1 & 4 \\ 13 & -4\end{array}\right| \\=& 7(-4+12)+1(-4-52) \\=& 56-56=0 \end{aligned}
\]

Singular. Rank <3, will need to put in echelon form to find exact rank. \\

\item[(d)]
\[ \begin{aligned} &-3\left|\begin{array}{ll}9 & 5 \\ 8 & 6\end{array}\right|-1\left|\begin{array}{cc}-4 & 9 \\ 10 & 8\end{array}\right| \\
 =&-3(54-40)-1(-32-90) \\
 =&-42+122=80 \end{aligned}\]
Non-singular, rank = 3. \vspace{1em}
\end{enumerate}

% Question 8
\item[8.] 
\begin{enumerate}
\item False. While we can find the transpose of any matrix, the determinant is only defined for square matrices. 
\item False. Multiplying each element of an $n \times n$ by 2 will increase the determinant by $2^n$ times. 
\item $A$ cannot be 0, but the determinant of $A$ can. In any case, if $|A|$=0, $A$ is singular.
\end{enumerate}

\end{enumerate}

%%%%%%%%%%%%%%%% Exercise 5.4

\underline{Exercise 5.4} 

\begin{enumerate}

% Question 2
\item[2.] Note that \[
A^{-1}=\frac{1}{|A|} \operatorname{adj} A
\]

Here, \( \operatorname{adj} A=C^{\prime} \) where \( C=\left[\left|C_{i j}\right|\right] \).

\begin{enumerate}

\item 
\[ |A|=5-0=5 \]
\[\operatorname{adj} A=\left[\begin{array}{ll}\left|C_{11}\right| & \left|C_{21}\right| \\ \left|C_{12}\right| & \left|C_{22}\right|\end{array}\right]=\left[\begin{array}{cc}1 & -2 \\ 0 & 5\end{array}\right] \]
\[A^{-1}=\frac{1}{|A|} \operatorname{adj} A=\frac{1}{5}\left[\begin{array}{cc}1 & -2 \\ 0 & 5\end{array}\right]=\left[\begin{array}{cc}1 / 5 & -2 / 5 \\ 0 & 1\end{array}\right] \]

\item \[
\begin{array}{l}
|B|=-2-0=-2 \\
B^{-1}=\frac{-1}{2}\left[\begin{array}{rr}
2 & 0 \\
-9 & -1
\end{array}\right]
\end{array} 
\]

\item \[
\begin{array}{l}
|C|=-3-21=-24 \\
C^{-1}=\frac{-1}{24}\left[\begin{array}{cc}
-1 & -7 \\
-3 & 3
\end{array}\right]
\end{array}
\]

\item \[
\begin{array}{l}
|D|=21-0=21 \\
D^{-1}=\frac{1}{21}\left[\begin{array}{rr}
3 & -6 \\
0 & 7
\end{array}\right]
\end{array}
\]
\end{enumerate}

% Question 3
\item[3.] 
\begin{enumerate}
\item Step 1: Exchange the two diagonal elements.\\
Step 2: Multiply both the off-diagonal elements by -1.

\item Step 3: Multiply the resulting matrix from steps 1 and 2 by \(1/|A|\).
\end{enumerate}

% Question 4
\item[4.] (a)\[
\begin{aligned}
|E| &=1\left|\begin{array}{ll}
7 & 3 \\
2 & 0
\end{array}\right|+1\left|\begin{array}{rr}
4 & -2 \\
7 & 3
\end{array}\right| \\
&=1(0-6)+1(12+14) \\
&=-6+26=20
\end{aligned}
\]

\[
\begin{aligned}
\operatorname{adj} E &=\left[\begin{array}{ccc}
\left|C_{11}\right| & \left|C_{21}\right| & \left|C_{31}\right| \\
\left|C_{12}\right| & \left|C_{22}\right| & \left|C_{32}\right| \\
\left|C_{13}\right| & \left|C_{23}\right| & \left|C_{33}\right|
\end{array}\right] \\
&=\left[\begin{array}{ccc}
3 & 2 & -3 \\
-7 & 2 & 7 \\
-6 & -4 & 26
\end{array}\right]
\end{aligned}
\]
\[ E^{-1}=\frac{1}{|E|} a d j E \] 

(b)
\[
F^{-1}=\frac{-1}{10}\left[\begin{array}{ccc}
0 & 2 & -3 \\
10 & -6 & -1 \\
0 & -4 & 1
\end{array}\right]
\]

(c)
\[ \left|C_{11}\right|=\left|\begin{array}{ll}0 & 1 \\ 1 & 0\end{array}\right|=-1 
\quad  \left|C_{12}\right|=-\left|\begin{array}{ll}0 & 1 \\ 0 & 0\end{array}\right|=0 
\quad  \left|C_{13}\right|=\left|\begin{array}{ll}0 & 0 \\ 0 & 1\end{array}\right|=0 
\]
\[ \left|C_{21}\right|=-\left|\begin{array}{ll}0 & 0 \\ 1 & 0\end{array}\right|=0 
\quad  \left|C_{22}\right|=\left|\begin{array}{ll}0 & 1 \\ 0 & 0\end{array}\right|=0 
\quad  \left|C_{23}\right|=-\left|\begin{array}{ll}1 & 0 \\ 0 & 1\end{array}\right|=-1 \]
\[
\left|C_{31}\right|=\left|\begin{array}{cc}
0 & 0 \\
0 & 1
\end{array}\right|=0 \quad\left|C_{32}\right|=-\left|\begin{array}{cc}
1 & 0 \\
0 & 1
\end{array}\right|=-1 \quad\left|C_{33}\right|=\left|\begin{array}{ll}
1 & 0 \\
0 & 0
\end{array}\right|=0
\]
\[ G^{-1}=-1\left[\begin{array}{ccc}-1 & 0 & 0 \\ 0 & 0 & -1 \\ 0 & -1 & 0\end{array}\right]=\left[\begin{array}{lll}1 & 0 & 0 \\ 0 & 0 & 1 \\ 0 & 1 & 0\end{array}\right] \] \\

(d) \[ H^{-1}=\left[\begin{array}{lll}1 & 0 & 0 \\ 0 & 1 & 0 \\ 0 & 0 & 1\end{array}\right] \] \\

% Question 7
\item[7.] If a matrix is its own inverse, we would need that \(A^2 = I \). This is true for the identity matrix. There are other possibilities such as matrix \(G\) in exercise 4. \vspace{1em}
\end{enumerate}
